\documentclass[aspectratio=169,10pt]{beamer}

\usetheme{default}
\usefonttheme{serif}
\setbeamertemplate{navigation symbols}{}
\setbeamertemplate{headline}{}

% Colors
\definecolor{cream}{HTML}{F5F3EE}
\definecolor{creamborder}{HTML}{D8D4CB}
\definecolor{darktext}{HTML}{1A1A18}
\definecolor{mutedtext}{HTML}{6B6860}

\setbeamercolor{background canvas}{bg=cream}
\setbeamercolor{normal text}{fg=darktext,bg=cream}
\setbeamercolor{frametitle}{fg=darktext,bg=cream}
\setbeamercolor{title}{fg=darktext,bg=cream}
\setbeamercolor{subtitle}{fg=mutedtext,bg=cream}
\setbeamercolor{author}{fg=darktext}
\setbeamercolor{date}{fg=mutedtext}
\setbeamercolor{block title}{fg=darktext,bg=creamborder}
\setbeamercolor{block body}{fg=darktext,bg=cream}
\setbeamercolor{block title alerted}{fg=darktext,bg=creamborder!70}
\setbeamercolor{block body alerted}{fg=darktext,bg=cream}
\setbeamercolor{alerted text}{fg=darktext}
\setbeamercolor{itemize item}{fg=darktext}
\setbeamercolor{itemize subitem}{fg=mutedtext}
\setbeamercolor{enumerate item}{fg=darktext}
\setbeamercolor{section in toc}{fg=darktext}

\usepackage[utf8]{inputenc}
\usepackage[T1]{fontenc}
\usepackage[english]{babel}
\usepackage{amsmath,amssymb}
\usepackage{graphicx}
\usepackage{booktabs}
\usepackage{multirow}
\usepackage{subcaption}
\usepackage{ebgaramond}

\graphicspath{{../outputs/}}

% Force correct landscape page dimensions in the output PDF
\AtBeginDocument{%
  \pdfpagewidth=\paperwidth%
  \pdfpageheight=\paperheight%
}

\setbeamerfont{frametitle}{size=\large,series=\bfseries}
\setbeamerfont{block title}{size=\small,series=\bfseries}

\setbeamertemplate{frametitle}{%
  \vspace{7pt}%
  \insertframetitle\par
  \vspace{3pt}%
  {\color{creamborder}\hrule height 0.5pt}%
  \vspace{4pt}%
}

\setbeamertemplate{footline}{%
  \hfill{\scriptsize\color{mutedtext}\insertframenumber}\hspace{10pt}\vspace{6pt}%
}

\setbeamertemplate{title page}{%
  \begin{beamercolorbox}[wd=\paperwidth,ht=\paperheight,dp=0pt]{background canvas}%
    \vspace{18pt}
    \hspace{28pt}%
    \begin{minipage}{0.82\paperwidth}
      {\fontsize{22}{28}\rmfamily\bfseries\inserttitle\par}
      \vspace{8pt}
      {\color{creamborder}\hrule height 0.6pt}
      \vspace{8pt}
      {\fontsize{8.5}{11}\rmfamily\color{mutedtext}\textsc{\insertsubtitle}\par}
      \vspace{4pt}
      {\fontsize{8.5}{11}\rmfamily\color{mutedtext}\textsc{\insertinstitute}\par}
    \end{minipage}
    \vfill
    \hspace{28pt}%
    \begin{minipage}{0.82\paperwidth}
      {\color{creamborder}\hrule height 0.4pt}
      \vspace{6pt}
      {\fontsize{9.5}{12}\rmfamily\insertauthor\quad{\color{mutedtext}|}\quad%
       {\color{mutedtext}\insertdate}\par}
    \end{minipage}
    \vspace{28pt}
  \end{beamercolorbox}%
}

\title{Generative Models\\for Cat Image Synthesis}
\subtitle{Deep Learning: Project III}
\author{Anna Ostrowska}
\institute{Warsaw University of Technology}
\date{June 2026}

\begin{document}

\begin{frame}[plain]
  \titlepage
\end{frame}

%
\section{Motivation \& Dataset}
%
\begin{frame}{Goal \& Motivation}
  \begin{columns}[T]
    \begin{column}{0.62\textwidth}
      Train and compare three generative architectures on cat images at
      $64\times64$ resolution with a 6\,GiB VRAM constraint.

      \vspace{6pt}
      \textbf{Three approaches}
      \begin{enumerate}
        \item \textbf{DCGAN} with Wasserstein loss and gradient penalty (WGAN-GP)
        \item \textbf{VQ-VAE} with autoregressive PixelCNN prior
        \item \textbf{DDPM} with DDIM fast sampling
      \end{enumerate}

      \vspace{6pt}
      \textbf{Evaluation:} Fréchet Inception Distance (FID) on 5\,000 samples,
      visual inspection, and latent interpolation.
    \end{column}
    \begin{column}{0.35\textwidth}
      \includegraphics[width=\linewidth]{data_samples_64.png}
      \centering\footnotesize Real cat samples ($64\times64$)
    \end{column}
  \end{columns}
\end{frame}

\begin{frame}{Dataset \& Preprocessing}
  \begin{columns}[T]
    \begin{column}{0.55\textwidth}
      \textbf{CAT Dataset (Kaggle)}
      \begin{itemize}
        \item 9\,997 RGB cat photos, varying resolutions
        \item Subdirectories CAT\_00 to CAT\_06
      \end{itemize}

      \vspace{6pt}
      \textbf{Preprocessing}
      \begin{enumerate}
        \item Resize shorter edge to 64\,px
        \item CenterCrop $64\times64$
        \item Normalize to $[-1,\,1]$
      \end{enumerate}

      \vspace{6pt}
      Fixed seed 42. No augmentation in baselines.
      Augmentation tested in ablation: FID $284.05\pm14.80$, worse than baseline $264.25\pm20.21$.
    \end{column}
    \begin{column}{0.42\textwidth}
      \includegraphics[width=\linewidth]{data_samples_64.png}
    \end{column}
  \end{columns}
\end{frame}

%
\section{Models}
%
\begin{frame}[shrink=8]{DCGAN with WGAN-GP}
  \begin{columns}[T]
    \begin{column}{0.54\textwidth}
      \textbf{Architecture}
      \begin{itemize}
        \item Generator: $\mathbf{z}\in\mathbb{R}^{128}$, 4 transposed-conv blocks, Tanh
        \item Discriminator: 4 strided-conv (LeakyReLU), scalar logit
      \end{itemize}

      \vspace{5pt}
      \textbf{Training}
      \begin{itemize}
        \item WGAN-GP: $\lambda_{GP}=10$, $n_\text{critic}=5$
        \item Adam $\beta_1=0.5$, lr $= 2\times10^{-4}$
        \item 100 epochs, batch 64
      \end{itemize}
      \vspace{3pt}
      \textbf{Note:} WGAN-GP critic outputs a scalar score, not a
      probability --- binary label targets have no role in this loss.

      \vspace{5pt}
      \textbf{Why WGAN-GP?} Smoother gradients, no weight-clipping issues.
      $n_\text{critic}=5$ keeps the critic well-calibrated before each
      generator update.
    \end{column}
    \begin{column}{0.43\textwidth}
      \includegraphics[width=\linewidth]{dcgan/samples_epoch0100.png}
      \centering\footnotesize Epoch 100
    \end{column}
  \end{columns}
\end{frame}

\begin{frame}[shrink=8]{VQ-VAE with PixelCNN Prior}
  \begin{columns}[T]
    \begin{column}{0.54\textwidth}
      \textbf{Phase 1: Autoencoder (100 epochs)}
      \begin{itemize}
        \item $64\times64$ to $16\times16$ discrete code map
        \item Codebook: $K=512$, $d=64$, EMA updates ($\gamma=0.99$)
        \item Commitment cost $\beta=0.25$
      \end{itemize}

      \vspace{5pt}
      \textbf{Phase 2: PixelCNN Prior (100 epochs)}
      \begin{itemize}
        \item 15-layer masked residual convolutions, 128 filters
        \item Trained on 256-token sequences (frozen encoder)
        \item Codebook perplexity: $403/512$ (79\% utilisation)
      \end{itemize}

      \vspace{5pt}
      \textbf{Why EMA codebook?} Prevents codebook collapse.
      Gradient-based updates are unstable; early run with $K=1024$
      had only 30\% code utilisation.
    \end{column}
    \begin{column}{0.43\textwidth}
      \includegraphics[width=\linewidth]{vqvae/prior_samples.png}
      \centering\footnotesize PixelCNN prior samples
    \end{column}
  \end{columns}
\end{frame}

\begin{frame}[shrink=8]{DDPM with DDIM Sampling}
  \begin{columns}[T]
    \begin{column}{0.54\textwidth}
      \textbf{Architecture}
      \begin{itemize}
        \item U-Net, channel mults $(1,2,2,2)$, 2 res-blocks per level
        \item Self-attention at $16\times16$; EMA weights at inference
      \end{itemize}

      \vspace{5pt}
      \textbf{Diffusion process}
      \begin{itemize}
        \item Linear schedule: $T=1000$, $\beta_1=10^{-4}$, $\beta_T=0.02$
        \item DDIM ($\eta=0$): 50 steps gives $20\times$ speedup
        \item 200 epochs, batch 16 (6 GiB VRAM limit)
      \end{itemize}

      \vspace{5pt}
      \textbf{Why EMA decay 0.9999?} Slower EMA reduces weight noise
      amplified by the long reverse chain. FID drops from 55.43 (no EMA)
      to 24.11 (with EMA) at 50 DDIM steps, a 57\% improvement.
    \end{column}
    \begin{column}{0.43\textwidth}
      \includegraphics[width=\linewidth]{ddpm/samples_epoch0200.png}
      \centering\footnotesize Epoch 200 (EMA weights)
    \end{column}
  \end{columns}
\end{frame}

%
\section{Results}
%
\begin{frame}[shrink=5]{FID Comparison}
  \begin{columns}[T]
    \begin{column}{0.52\textwidth}
      \begin{block}{Fréchet Inception Distance (5\,000 samples)}
        \centering\vspace{3pt}
        \begin{tabular}{lcc}
          \toprule
          Model & FID $\downarrow$ & Training \\
          \midrule
          DCGAN (WGAN-GP, $n_c{=}5$) & 221.51 & $\sim$2\,h \\
          \quad$\hookrightarrow$ corrected (std, $n_c{=}1$) & 98.99 & $\sim$3\,h \\
          VQ-VAE+PixelCNN   & 187.54 & $\sim$4\,h \\
          \textbf{DDPM+EMA} & \textbf{24.11}  & $\sim$12\,h \\
          \bottomrule
        \end{tabular}\vspace{3pt}
      \end{block}
      \vspace{2pt}
      \footnotesize Not compute-normalised; shows best feasible config per model.
      All single-run. Original DCGAN was under-trained (see §DCGAN Fix).
    \end{column}
    \begin{column}{0.46\textwidth}
      \textbf{Why such a large gap?}
      \begin{itemize}
        \item DDPM has a well-defined MSE objective: smooth, stable gradients
        \item EMA weights reduce FID by 57\% vs.\ non-EMA at 50 steps
        \item GAN mode collapse limits diversity despite WGAN-GP
        \item VQ-VAE: good reconstruction, PixelCNN prior limits sample quality
      \end{itemize}
      \vspace{4pt}
      \footnotesize Consistent with \textit{Diffusion Models Beat GANs} (Dhariwal \& Nichol, 2021)
    \end{column}
  \end{columns}
\end{frame}

\begin{frame}{Generated Samples: Side by Side}
  \begin{figure}
    \centering
    \begin{subfigure}[b]{0.31\textwidth}
      \includegraphics[width=\linewidth]{dcgan/samples_epoch0100.png}
      \caption{DCGAN\\{\small FID 221.51}}
    \end{subfigure}
    \hfill
    \begin{subfigure}[b]{0.31\textwidth}
      \includegraphics[width=\linewidth]{vqvae/prior_samples.png}
      \caption{VQ-VAE+PixelCNN\\{\small FID 187.54}}
    \end{subfigure}
    \hfill
    \begin{subfigure}[b]{0.31\textwidth}
      \includegraphics[width=\linewidth]{ddpm/samples_epoch0200.png}
      \caption{\textbf{DDPM}\\{\small FID 24.11}}
    \end{subfigure}
  \end{figure}
  \vspace{1pt}
  \small
  DCGAN: artifacts, limited diversity \quad|\quad
  VQ-VAE: plausible but blurry \quad|\quad
  \textbf{DDPM: sharpest, most coherent}
\end{frame}

\begin{frame}{Real vs.\ Generated}
  \centering
  \includegraphics[width=\textwidth,height=0.82\textheight,keepaspectratio]{real_vs_generated.png}
\end{frame}


\begin{frame}[shrink=8]{Latent Space Interpolation: Linear vs.\ Slerp}
  \centering
  \includegraphics[width=\textwidth,height=0.28\textheight,keepaspectratio]{dcgan/interpolation_compare.png}\\
  \small DCGAN --- \textbf{linear} (top) vs.\ slerp (bottom)\\[5pt]
  \includegraphics[width=\textwidth,height=0.28\textheight,keepaspectratio]{ddpm/interpolation_compare.png}\\
  \small DDPM --- \textbf{linear} (top) vs.\ slerp (bottom)\\[4pt]
  \footnotesize\color{mutedtext}
  Linear interpolation is the primary method (top rows).
  Slerp is shown for comparison: it preserves the norm of the latent vector,
  avoiding the low-density region near the origin of $\mathcal{N}(0,I)$
  that causes blurry midpoints in linear interpolation (most visible in DDPM).
\end{frame}

\begin{frame}[shrink=8]{VQ-VAE Reconstruction Quality}
  \begin{center}
    \includegraphics[width=0.90\linewidth]{vqvae/recon_slide.png}
  \end{center}
  \vspace{1pt}
  \small\centering
  Top: input \quad|\quad Bottom: VQ-VAE reconstruction (epoch 100)\\[1pt]
  MSE $= 0.0066$ \quad|\quad VQ loss $= 0.0033$ \quad|\quad perplexity $403/512$ (79\%)
\end{frame}


%
\section{Ablation Studies}
%
\begin{frame}{DDPM: DDIM Steps and EMA Effect}
\small
  \begin{columns}[T]
    \begin{column}{0.48\textwidth}
      \begin{block}{FID vs. DDIM steps (non-EMA)}
        \centering\vspace{3pt}
        \begin{tabular}{ccc}
          \toprule
          Steps & FID $\downarrow$ & Speedup \\
          \midrule
          10  & 66.90 & $100\times$ \\
          25  & 60.26 & $40\times$ \\
          \textbf{50}  & \textbf{55.43} & $\mathbf{20\times}$ \\
          100 & 48.35 & $10\times$ \\
          200 & 39.56 & $5\times$ \\
          \bottomrule
        \end{tabular}\vspace{3pt}
      \end{block}
    \end{column}
    \begin{column}{0.48\textwidth}
      \begin{block}{EMA at 50 steps: $-57\%$ FID}
        \centering\vspace{3pt}
        \begin{tabular}{lc}
          \toprule
          Weights & FID $\downarrow$ \\
          \midrule
          Without EMA & 55.43 \\
          \textbf{With EMA} & \textbf{24.11} \\
          \bottomrule
        \end{tabular}\vspace{3pt}
      \end{block}
      \vspace{4pt}
      50 steps: practical sweet spot ($20\times$ speedup).\\
      EMA more impactful than doubling the steps.\\
      Below 25 steps: FID degrades by 37\% vs. 100 steps.
    \end{column}
  \end{columns}
\end{frame}

\begin{frame}[shrink=5]{DCGAN: Latent Dimension and Learning Rate}
  \begin{columns}[T]
    \begin{column}{0.48\textwidth}
      \begin{block}{$z_\text{dim}$ ablation (3 seeds)}
        \centering\vspace{3pt}
        \begin{tabular}{lcc}
          \toprule
          $z_\text{dim}$ & FID mean$\pm$std & $\Delta$ \\
          \midrule
          64  & $246.65 \pm 5.57$  & $-17.60$ \\
          128 (baseline) & $264.25 \pm 20.21$ & - \\
          256 & $266.05 \pm 2.52$  & $+1.80$ \\
          \bottomrule
        \end{tabular}\vspace{3pt}
      \end{block}
      \vspace{3pt}
      $z_{64}$ best, but within baseline $\sigma$. High baseline variance ($\pm20$) shows seed sensitivity.
    \end{column}
    \begin{column}{0.48\textwidth}
      \begin{block}{Learning rate ablation (3 seeds)}
        \centering\vspace{3pt}
        \begin{tabular}{lcc}
          \toprule
          LR & FID mean$\pm$std & $\Delta$ \\
          \midrule
          $1\times10^{-4}$ & $295.11 \pm 18.27$ & $+30.86$ \\
          $2\times10^{-4}$ (base) & $264.25 \pm 20.21$ & - \\
          \textbf{$4\times10^{-4}$} & $\mathbf{220.05 \pm 4.47}$ & $-44.20$ \\
          \bottomrule
        \end{tabular}\vspace{3pt}
      \end{block}
      \vspace{3pt}
      LR has far more impact than $z_\text{dim}$.\\
      Low LR: generator cannot keep up with the critic.
    \end{column}
  \end{columns}
\end{frame}


%
\section{Cats and Dogs Experiment}
%
\begin{frame}[shrink=18]{Does Training on Dogs Help Cats?}
  \begin{columns}[T]
    \begin{column}{0.52\textwidth}
      \textbf{Setup}
      \begin{itemize}
        \item 12\,500 cats + 12\,500 dogs (Kaggle)
        \item Same DCGAN, 100 epochs, unconditional
        \item FID measured separately vs. cats and dogs
      \end{itemize}

      \vspace{4pt}
      \begin{block}{Results}
        \centering\vspace{2pt}
        \begin{tabular}{lccc}
          \toprule
          Model & vs. Cats & vs. Dogs & vs. Combined \\
          \midrule
          Cats only  & 221.51 & - & - \\
          \textbf{Cats+Dogs} & 206.57 & 197.79 & \textbf{190.39} \\
          \bottomrule
        \end{tabular}\vspace{2pt}
      \end{block}

      \vspace{3pt}
      \small Combined FID (2\,500 cats + 2\,500 dogs) more appropriate
      for a mixed-class model. 15-pt cat improvement is within
      seed variance ($\pm20$) --- suggestive, not conclusive.

      \vspace{3pt}
      \textbf{Qualitative:} mix of cat-like and dog-like samples;
      no collapse to one class. Model captures shared fur/eye features
      without cleanly separating species.
    \end{column}
    \begin{column}{0.45\textwidth}
      \includegraphics[width=\linewidth]{cats_and_dogs/samples_epoch0100.png}
      \centering\footnotesize Cats+Dogs DCGAN (epoch 100)
    \end{column}
  \end{columns}
\end{frame}

%
\section{DCGAN: Training Bottleneck \& Fix}
%
\begin{frame}[shrink=10]{DCGAN: Why FID Was So High and How It Was Fixed}
\small
  \begin{columns}[T]
    \begin{column}{0.52\textwidth}
      \textbf{Post-hoc diagnosis: generator update budget}
      \begin{itemize}
        \item $n_\text{critic}=5$: 5 D-updates per G-update
        \item 156 batches/epoch $\Rightarrow$ only \textbf{3\,100 G-steps} in 100 ep.
        \item Standard DCGAN ($n_c=1$): \textbf{15\,600} at same cost
      \end{itemize}

      \vspace{4pt}
      \textbf{Two corrected configs}
      \begin{itemize}
        \item \textbf{Standard DCGAN}: switch to BCE loss; $n_c=1$ is the
              correct standard (not a compromise), $g_\text{ch}=128$,
              150 ep.\ $\Rightarrow$ \textbf{23\,400 G-steps}
        \item \textbf{WGAN-GP $n_c=2$}: pragmatic fix; fewer critic steps
              may leave the critic less well-optimised ---
              compute-equivalent $n_c{=}5$ would need $\approx$500 ep., 200 ep.\ $\Rightarrow$ \textbf{15\,600 G-steps}
      \end{itemize}

      \vspace{4pt}
      \begin{block}{FID results}
        \centering\vspace{2pt}
        \begin{tabular}{lcc}
          \toprule
          Config & G-steps & FID \\
          \midrule
          Original       & 3\,100  & 221.51 \\
          Std.\ DCGAN    & 23\,400 & \textbf{98.99} \\
          WGAN-GP $n_c=2$ & 15\,600 & 158.48 \\
          \bottomrule
        \end{tabular}\vspace{2pt}
      \end{block}
    \end{column}
    \begin{column}{0.45\textwidth}
      \centering
      \includegraphics[width=\linewidth,height=0.38\textheight,keepaspectratio]{dcgan_fixed_vanilla/samples_epoch0150.png}\\[2pt]
      \footnotesize Std.\ DCGAN (epoch 150)\\[6pt]
      \includegraphics[width=\linewidth,height=0.38\textheight,keepaspectratio]{dcgan_fixed_wgangp/samples_epoch0200.png}\\[2pt]
      \footnotesize WGAN-GP $n_c=2$ (epoch 200)
    \end{column}
  \end{columns}
\end{frame}


%
\section{Research Questions}
%
\begin{frame}[shrink=5]{Research Questions: Answers}
\footnotesize
  \begin{block}{RQ1: Architecture trade-off under 6\,GiB VRAM}
    DDPM: best quality (FID 24.11, 12\,h, 5.5\,GiB).
    DCGAN: fastest (2\,h, $<$2\,GiB) but mode collapse.
    VQ-VAE: stable training, PixelCNN prior is the bottleneck.
  \end{block}
  \vspace{3pt}
  \begin{block}{RQ2: Does adding dogs improve cat generation?}
    Possibly, but not conclusively. Cat FID improves from 221.51 to 206.57, but the
    15-pt gap is within the observed DCGAN seed variance ($\pm20$). Class coverage
    and absence of hybrid artefacts would require a classifier-based analysis.
  \end{block}
  \vspace{3pt}
  \begin{block}{RQ3: Effect of latent dimension $z_\text{dim} \in \{64,128,256\}$}
    Differences are within baseline seed variance ($\pm20$); no firm conclusion about
    optimal $z_\text{dim}$ can be drawn. $z=64$ best on mean FID but not conclusively.
  \end{block}
  \vspace{3pt}
  \begin{block}{RQ4: DDIM speed-quality trade-off}
    50 steps: $20\times$ speedup, FID 55.43 (non-EMA) or 24.11 (with EMA).
    Below 25 steps quality degrades sharply. 50 steps is the practical sweet spot.
  \end{block}
\end{frame}

%
\section{Conclusion}
%
\begin{frame}[shrink=8]{Conclusion}
\small
  \begin{columns}[T]
    \begin{column}{0.52\textwidth}
      \begin{block}{Summary}
        \centering\vspace{3pt}
        \begin{tabular}{lcc}
          \toprule
          Model & FID & Strength \\
          \midrule
          \textbf{DDPM} & \textbf{24.11} & quality \\
          VQ-VAE        & 187.54 & stability \\
          DCGAN         & 221.51 & speed \\
          \bottomrule
        \end{tabular}\vspace{3pt}
      \end{block}
      \vspace{5pt}
      \begin{itemize}
        \item In our setup, DDPM clearly outperforms GAN and VAE
        \item For DDPM, EMA more impactful than doubling DDIM steps
        \item DDIM enables practical use at 50 steps
        \item Dog images: suggestive improvement, not conclusive
        \item DCGAN bottleneck fixed: $7\times$ more G-updates $\to$ FID 98.99
      \end{itemize}
    \end{column}
    \begin{column}{0.46\textwidth}
      \textbf{Main challenges}
      \begin{itemize}
        \item 6\,GiB VRAM: DDPM batch size limited to 16
        \item DCGAN mode collapse: WGAN-GP (label smoothing inactive under GP path)
        \item Codebook collapse at $K=1024$: fixed with EMA, $K=512$
        \item 8\,h inference: solved with DDIM (50 steps, 25\,min)
        \item DCGAN $n_\text{critic}=5$: generator starvation fixed
      \end{itemize}
      \vspace{6pt}
      \textbf{Future work}
      \begin{itemize}
        \item DDPM at $128\times128$ with gradient checkpointing
        \item Conditional DCGAN on cats+dogs
        \item Transformer prior (VQ-VAE-2 style)
        \item Cosine noise schedule
      \end{itemize}
    \end{column}
  \end{columns}
\end{frame}

\begin{frame}[plain]
  \vfill
  \begin{center}
    {\fontsize{34}{42}\rmfamily\bfseries Thank you.}
    \vspace{16pt}

    {\color{creamborder}\hrule height 0.5pt}
    \vspace{12pt}

    {\normalsize Anna Ostrowska}\\[4pt]
    {\small\color{mutedtext}%
      Deep Learning: Project III \quad|\quad
      Warsaw University of Technology \quad|\quad
      June 2026}
  \end{center}
  \vfill
\end{frame}


\appendix

\begin{frame}[shrink=12]{Training Dynamics}
  \begin{columns}[T]
    \begin{column}{0.48\textwidth}\centering
      \includegraphics[width=\linewidth,height=0.38\textheight,keepaspectratio]{dcgan/losses.png}\\[2pt]
      DCGAN (G vs.\ D) --- adversarial oscillation
    \end{column}
    \begin{column}{0.48\textwidth}\centering
      \includegraphics[width=\linewidth,height=0.38\textheight,keepaspectratio]{vqvae/losses.png}\\[2pt]
      VQ-VAE (recon + VQ) --- steady decrease
    \end{column}
  \end{columns}
  \vspace{4pt}
  \begin{center}
    \includegraphics[width=0.48\textwidth,height=0.38\textheight,keepaspectratio]{ddpm/losses.png}\\[2pt]
    DDPM ($\ell_2$ noise) --- smooth convergence, most stable
  \end{center}
  \vspace{3pt}
  \small DDPM's MSE objective gives the most stable training.
  DCGAN shows adversarial oscillation; VQ-VAE converges reliably with EMA codebook.
\end{frame}

\begin{frame}{DCGAN Training Progression}
  \centering
  \includegraphics[width=\textwidth,height=0.78\textheight,keepaspectratio]{dcgan_progression.png}\\[3pt]
  \small From blurry noise to recognisable cats --- but likely limited by low generator update budget (3\,100 G-steps at $n_c{=}5$).
\end{frame}

\begin{frame}[shrink=5]{Experimental Grid: What We Tested}
\small
  \begin{tabular}{llll}
    \toprule
    \textbf{Model} & \textbf{Ablation / variant} & \textbf{Values tested} & \textbf{Metric} \\
    \midrule
    \multirow{4}{*}{DCGAN}
      & Latent dim $z_\text{dim}$   & 64, \textbf{128}, 256              & FID, slerp smoothness \\
      & Learning rate               & 1e-4, \textbf{2e-4}, 4e-4          & FID \\
      & Augmentation                & off (\textbf{base}), on            & FID \\
      & Training data               & cats only, \textbf{cats+dogs}       & FID vs cats/dogs \\
    \midrule
    \multirow{2}{*}{DDPM}
      & DDIM steps                  & 10, 25, \textbf{50}, 100, 200      & FID, wall time \\
      & EMA weights                 & off, \textbf{on} (decay 0.9999)    & FID \\
    \midrule
    VQ-VAE
      & Codebook size $K$           & 1024 (collapsed) $\to$ \textbf{512}   & perplexity, FID \\
    \midrule
    \multirow{2}{*}{Post-hoc fix}
      & Standard DCGAN ($n_c=1$)    & 150 ep, $g_\text{ch}=128$          & FID \\
      & WGAN-GP $n_c=2$             & 200 ep                              & FID \\
    \bottomrule
  \end{tabular}
  \vspace{5pt}

  \textbf{Bold} = chosen baseline.
  Primary metric: FID on 5\,000 generated vs 5\,000 real images.
  Secondary: latent interpolation smoothness, reconstruction MSE (VQ-VAE).
\end{frame}

\begin{frame}[shrink=5]{Latent Space Smoothness vs. $z_\text{dim}$}
  \centering
  \includegraphics[width=\textwidth,height=0.20\textheight,keepaspectratio]{ablation/interp_zdim64.png}\\
  \small $z_\text{dim}=64$\\[3pt]
  \includegraphics[width=\textwidth,height=0.20\textheight,keepaspectratio]{ablation/interp_zdim128.png}\\
  \small $z_\text{dim}=128$ (baseline)\\[3pt]
  \includegraphics[width=\textwidth,height=0.20\textheight,keepaspectratio]{ablation/interp_zdim256.png}\\
  \small $z_\text{dim}=256$\\[4pt]
  Slerp, 10 frames. Transitions smooth for all variants.
  $z_{64}$ slightly softer; $z_{128}$ and $z_{256}$ indistinguishable.
\end{frame}

\begin{frame}[shrink=5]{Quality vs.\ Compute Trade-off}
  \centering
  \includegraphics[width=0.88\textwidth,height=0.80\textheight,keepaspectratio]{compute_quality.png}\\[2pt]
  \small DDPM dominates in quality; DCGAN fastest. Fixed DCGAN closes the gap significantly.
\end{frame}

\end{document}
